\begin{prob}
    While on campus, you notice someone standing next to a very tall ladder
    getting ready to replace a lightbulb. As they are about to get started,
    though, they are distracted by a raccoon and accidentally drop the bulb to
    the ground. Surprisingly, though, the bulb doesn't break!

    You wonder to yourself: how high up could the bulb be dropped without it
    breaking? The person goes away for a moment, leaving their ladder, two
    bulbs, and an opportunity for you to find out the answer to your question.
    They'll be back soon, though, so you must hurry. You decide to test the
    bulbs by dropping them and seeing when they break.

    More formally, the ladder has $n$ rungs (that is, $n$ steps). You want to
    find out the maximum step, $n_\text{max}$, that you can drop a bulb without
    it breaking. You'll assume that the two bulbs are both equally-strong, and
    will break on the same step. Since time is limited, you want to find out the
    answer to your question in as few bulb-drops as you can. You're allowed to
    break both bulbs.

    One approach is essentially linear search. Here, you stand on step 1, drop
    the bulb, and see if it breaks. If it doesn't, move to step 2, and so on. If
    the bulb breaks on step $k$, then you know the highest you can drop the bulb
    from is step $k - 1$. While this strategy is guaranteed to find you the
    answer and breaks only one bulb, it is time consuming: in the worst case,
    you'll need to drop the bulb $\Theta(n)$ times.

    But you've taken DSC 40B, so you're clever -- what about binary search?
    In this approach, you'll start on step $n / 2$, and drop the first bulb --
    if it doesn't break, you'll go higher.
    This seems more efficient, but there's a big problem with this: what if you
    break the bulb on the first drop? Then you only have one bulb remaining,
    and you have to be careful. You'll need to go back to using linear search,
    dropping the remaining bulb from step 1, step 2, and so forth until it
    breaks. In the worst case this linear search will do around $n/2 =
    \Theta(n)$ drops.

    However, there is a strategy (many strategies, actually) for finding out
    $n_\text{max}$ by breaking at most two bulbs and in the worst case making a
    number of drops $f(n)$ that is asymptotically much smaller than $n$. That
    is, if $f(n)$ is the number of drops needed by your method in the worst
    case, it should be true that $\lim_{n \to \infty} f(n) / n = 0$.

    Give such a strategy and state the number of drops it needs in the worst
    case, $f(n)$, using asymptotic notation. There are many strategies that
    satisfy the conditions -- yours does not need to be the most efficient.
    
    \begin{soln}
        Go up to step $\sqrt{n}$ and drop the first bulb. If it doesn't break,
        go up to step $2\sqrt{n}$ and drop. Repeat until the bulb breaks. Say
        it breaks at step $k\sqrt{n}$. Then you only need to check steps
        $(k-1)\sqrt n$ through $k \sqrt n$, which you can do with a linear
        search in less than $\sqrt{n}$ steps.

        In the worst case, $n_\text{max} = n-1$. In this case, you'll drop
        at $\sqrt n$, then $2 \sqrt n$, then $3 \sqrt n$ and so on, all the way
        to the top. It takes $n / \sqrt{n} = \sqrt n$ drops to get to the top.
        Then the bulb finally breaks, and you have to do a linear search to
        find $n_\text{max}$ exactly -- it could be anywhere between the top
        rung, $n$, and the location of the previous drop, $n - \sqrt{n}$.
        This linear search takes about $\sqrt n$ drops.
        The total number of drops is $\sqrt n$.

        A concrete example might make this clearer. Say $n = 100$. Then
        start on step 10 and drop the bulb. Say it doesn't break. Move to step
        20 and drop the first bulb again. Still doesn't break. Now
        drop from 30 -- supposed it breaks. $n_\text{max}$ is between 20 and 30,
        so drop the remaining bulb from 21, 22, 23, etc. until it breaks. In the
        worst case, you drop the bulb from $10, 20, 30, \ldots, 90, 100$, and it
        breaks at 100, so you then drop from $91, 92, \ldots, 99$.

        Note that $\sqrt n$ is (close) to the optimal solution, but there are 
        infinitely many others. For example, you could drop from $\log n$, then $2 \log n$, etc.
        Here the worst case would be $\Theta(n / \log n)$, which is still better
        than $\Theta(n)$. Similarly, you could drop at $n^{1/4}$, etc.

        The most efficient solution is determined by the number of bulbs you
        have. If you have $\Theta(\log n)$ bulbs, the most efficient solution is
        binary search.

    \end{soln}

\end{prob}
